\subsection{Survey Scope}

本研究のsurveyは、量子VRP/CVRP/EVRP研究を網羅的に列挙することではなく、stage-based gap assessmentに必要な証拠を抽出できる文献群を整理することを目的とする。
そのため、survey対象は三つの役割に分ける。
第一に、QAOA、NISQ、VQA、application-oriented benchmark、ARL、resource estimationなど、分析枠組みを支える概念・方法論文献である。
第二に、Figure~\ref{fig:width}とTable~\ref{tab:gap-matrix}を作成するために、width、depth、gate count、shots、feasibility、workflow情報を抽出した分析対象文献である。
第三に、社会側readinessやEVRP制約、dynamic routing、量子技術ロードマップを位置づける背景資料である。

\subsection{Survey Papers to Be Covered}

Table~\ref{tab:survey-paper-list} lists the papers and source documents that should be covered in the survey section.
The table separates sources used as conceptual or methodological references from papers used as the extraction corpus.
Papers in the extraction corpus are listed here to clarify the survey scope, but they are not cited through BibTeX in the main text; their bibliographic details are instead provided in Appendix~\ref{app:analysis-corpus}.

\begingroup
\scriptsize
\setlength{\tabcolsep}{2pt}
\begin{longtable}{>{\raggedright\arraybackslash}p{0.15\linewidth}>{\raggedright\arraybackslash}p{0.30\linewidth}>{\raggedright\arraybackslash}p{0.17\linewidth}>{\raggedright\arraybackslash}p{0.27\linewidth}}
\caption{Papers and source documents to be covered in the survey.}
\label{tab:survey-paper-list}\\
\toprule
Source ID & Paper or source & Survey role & What to extract or discuss \\
\midrule
\endfirsthead
\toprule
Source ID & Paper or source & Survey role & What to extract or discuss \\
\midrule
\endhead
qaoa-2014 & Farhi et al., ``A Quantum Approximate Optimization Algorithm'' & Conceptual method & QAOA layer structure, $p$, cost/mixer alternation, definition-level width/depth. \\
nisq-2018 & Preskill, ``Quantum Computing in the NISQ era and beyond'' & Conceptual method & NISQ noise, limited reliable circuit size, need for higher-fidelity gates and FTQC. \\
vqa-2021 & Cerezo et al., ``Variational quantum algorithms'' & Conceptual method & Hybrid quantum-classical workflow, optimizer loop, NISQ relevance and limitations. \\
bench-2021 & Lubinski et al., ``Application-oriented performance benchmarks for quantum computing'' & Benchmark method & Width/depth as application-oriented benchmark dimensions and volumetric evaluation. \\
bench-2023 & Lubinski et al., ``Optimization applications as quantum performance benchmarks'' & Benchmark method & Optimization benchmark design and interpretation of application-level quantum performance. \\
bench-2024 & Lubinski et al., ``Quantum algorithm exploration using application-oriented performance benchmarks'' & Benchmark method & Application-oriented benchmark extension and evidence classification. \\
utility-2023 & Herrmann et al., ``Quantum utility -- definition and assessment of a practical quantum advantage'' & Readiness method & Quantum utility, application readiness levels, and distinction from this study's QRL. \\
resource-2022 & Beverland et al., ``Assessing requirements to scale to practical quantum advantage'' & Resource-estimation method & Full-stack resource estimation and practical quantum advantage requirements. \\
vrp-2020 & Azad et al., ``Solving Vehicle Routing Problem Using Quantum Approximate Optimization Algorithm'' & Analysis corpus & Small VRP-QAOA simulation width/depth; early VRP-QAOA baseline. \\
vrp-2025 & Azfar et al., ``Quantum-Assisted Vehicle Routing'' & Analysis corpus & Gate-based VRP-QAOA, qubits, transpiled gates, shots/runtime, infeasible-state warning. \\
hvrp-2024 & Fitzek et al., ``Applying quantum approximate optimization to the heterogeneous vehicle routing problem'' & Analysis corpus & HVRP QAOA instances, qubits, $p$, gate-scaling discussion. \\
vrptw-2023 & Leonidas et al., ``Qubit efficient quantum algorithms for the vehicle routing problem on NISQ processors'' & Analysis corpus & Full/minimal encoding, route variables, qubits, ansatz layers, backend comparison. \\
cvrp-2023a & Palackal et al., ``Quantum-Assisted Solution Paths for the Capacitated Vehicle Routing Problem'' & Analysis corpus & CVRP-to-TSP decomposition, QAOA/VQE comparison, feasibility and runtime warnings. \\
cvrp-2023b & Xie et al., ``A Feasibility-Preserved Quantum Approximate Solver for the Capacitated Vehicle Routing Problem'' & Analysis corpus & Feasibility-preserving AOA, register-based width formula, depth 1--9 experiments, gate complexity. \\
cvrp-2025 & Onah and Michielsen, ``Requirements for Early Quantum Utility and Quantum Utility in the Capacitated Vehicle Routing Problem'' & Analysis corpus & CVRP resource estimation, HOBO/QUBO width, depth, quantum volume and utility gap. \\
encoding-2022 & Glos et al., ``Space-efficient binary optimization for variational quantum computing'' & Analysis corpus & QUBO/HOBO/mixed encoding width-depth trade-off for routing-style binary optimization. \\
logistics-2024 & Phillipson et al., logistics quantum-computing overview & Background context & Logistics quantum-computing landscape and current practical barriers. \\
dynamic-vrp-2013 & Pillac et al., ``A review of dynamic vehicle routing problems'' & Social context & Dynamic VRP concepts and operational complexity for SRL definition. \\
evrptw-2014 & Schneider et al., ``The electric vehicle-routing problem with time windows and recharging stations'' & Social context & EVRP constraints, time windows, charging stations, EV-specific routing complexity. \\
ev-routing-2020 & Frederhausen and Klumpp, commercial EV routing review & Social context & Urban logistics and commercial EV routing conditions. \\
doe-roadmap-2024 & U.S. DOE, ``Quantum Information Science Applications Roadmap'' & Roadmap context & Quantum technology roadmap and future-stage comparison. \\
trl-2020 & NASA technology readiness assessment guidance & Readiness method & General readiness-level logic and caution in interpreting stage labels. \\
systems-2000 & Sterman, ``Business Dynamics'' & Systems method & Systems-thinking framing; not used as full stock-flow simulation evidence. \\
socio-tech-2002 & Geels, socio-technical transition framework & Socio-technical method & Socio-technical transition framing for connecting social demand and technology readiness. \\
\bottomrule
\end{longtable}
\endgroup

\subsection{Conceptual and Methodological Sources}

概念・方法論文献は、量子アルゴリズムの構造、NISQ制約、hybrid workflow、アプリケーション志向ベンチマーク、量子utility、resource estimationの考え方を定義するために用いる。
QAOAの基本構造はFarhi et al.に基づき、NISQにおけるノイズと回路サイズ制約はPreskillに基づく。
VQA/QAOAのhybrid quantum-classical workflowはCerezo et al.を参照する。
Application-oriented performance benchmarkとwidth/depthを用いたvolumetricな評価視点はLubinski et al.を参照する。
Quantum utilityとapplication readinessの段階評価はHerrmann et al.を参照し、実用的量子優位に必要なfull-stack resource estimationの観点はBeverland et al.を参照する
\cite{farhi2014qaoa,preskill2018nisq,cerezo2021vqa,lubinski2021applicationbenchmarks,lubinski2024optimizationbenchmarks,lubinski2024applicationbenchmarks,herrmann2023quantumutility,beverland2022practicaladvantage}。

\subsection{Analysis Corpus for Circuit-Resource Extraction}

分析対象文献は、概念・方法論の引用文献とは区別する。
これらは一般的主張を支えるReferencesとしてではなく、図表作成のために調査し、回路資源・workflow証拠を抽出したcorpusとして扱う。
分析対象文献から抽出した具体データは、\texttt{circuit\_resources.csv} に記録し、Appendix~\ref{app:extracted-circuit-resources}に再掲する。
分析対象文献の一覧はAppendix~\ref{app:analysis-corpus}に記載する。

このcorpusでは、同一論文内でもinstance、encoding、実行環境が異なる場合は別の抽出単位として扱う。
例えば、full encodingとminimal encoding、QUBOとHOBO、small simulationとhardware-oriented execution、resource estimationは同じ意味のデータではない。
したがって、surveyでは単に「どの論文があるか」ではなく、「どの文献から、どの種類の証拠が抽出できるか」を整理する。

\subsection{Social and Roadmap Sources}

社会側readinessを定義するためには、EV stockを単一のVRP車両数へ直接変換するのではなく、EV・物流システムにおける制約条件、運用複雑性、dynamic routing、充電制約、都市物流条件を段階として整理する。
そのため、dynamic VRP、EVRPTW、commercial EV routing、量子技術ロードマップに関する資料を背景資料として用いる
\cite{pillac2013dynamicvrp,schneider2014evrptw,frederhausen2020commercialevrouting,doe2024qisapplicationsroadmap}。

これらの背景資料は、量子回路資源値を直接与えるものではない。
本研究では、社会側資料をSRLの定義・制約軸の整理に用い、量子側のwidth/depth抽出データとは別レイヤーとして扱う。
